\subsection{Community-Building: Rettungs-Feed und Challenges}\label{sec:community_und_feedback}

Der Rettungs-Feed ist der zentrale Ort, an dem Nutzer:innen ihre Reste-Kreationen teilen. Jeder Beitrag besteht aus einem Foto oder kurzen Video und einem Tag, der die verwendete Reste-Zutat markiert, zum Beispiel \enquote{Brokkoli-Strunk}. Andere können Beiträge kommentieren und bewerten, wodurch der Feed genau den Austausch ermöglicht, den die Wettbewerbsanalyse bei den reinen Reste-Matching-Apps vermisst hat. Durch Antippen des bei jedem Beitrag sichtbaren Nutzernamens lässt sich anderen zudem folgen.

Der Feed baut auf einem Mechanismus auf, den \textcite[S. 4, 6]{shenInfluenceOsmosisSocial2023} als \enquote{Influence by Osmosis} beschreiben, also Einfluss durch bloße Sichtbarkeit: Menschen übernehmen umweltfreundliches Verhalten häufig, weil sie es bei anderen sehen, nicht weil sie dazu aufgefordert werden. Für Resteria bedeutet das, dass bereits das Durchscrollen fremder Reste-Kreationen die eigene Motivation erhöhen kann, unabhängig von eigener Aktivität.

Resteria organisiert die Interaktion zwischen Nutzer:innen vor allem über zeitlich begrenzte Challenges, etwa eine monatliche \enquote{Zero-Waste-Woche} mit Rangliste, bei der die meisten geretteten Zutaten oder die kreativsten Reste-Gerichte prämiert werden. Erwogen, aber verworfen wurden dagegen ein offenes Diskussionsforum sowie feste thematische oder regionale Nutzergruppen. Den Anstoß dazu gab ein Befund von \textcite[S. 10--12]{somaFoodWasteReduction2020}. Dort erreichte eine Gruppe mit gemeinsamen Präsenzveranstaltungen deutlich geringere Teilnahme als eine passive Informationsgruppe oder eine Gamification-Gruppe. Forum und Nutzergruppen erfordern wie diese Präsenzgruppe aktiven Gruppenaustausch. Ob sich der Teilnahmeunterschied auf Online-Formate übertragen lässt, ist eine Annahme dieser Arbeit. Resteria ersetzt Forum und Gruppen deshalb durch zeitlich begrenzte Challenges. Der Rettungs-Feed entspricht dabei bereits der passiven Informationsgruppe aus der Studie, die dort ähnlich gut abschnitt wie die Gamification-Gruppe. Rettungspunkte und Reste-Level setzen den dort beobachteten Unterschied zwischen stärker und schwächer engagierten Nutzer:innen obendrauf (siehe~\autoref{sec:konzept}).
